This chapter fits one risk model per endpoint on the cohort of Chapter 4 —
five binary outcomes at 335 days, five models — and reports how well each of
them discriminates. The models are not an end in themselves. They exist for
two downstream uses and one incidental finding.

The first use is structural. The two-model meta-learner of Chapter 6 estimates
the individual effect as the difference between two outcome models, one fitted
per treatment arm; those two models are exactly the risk models of this
chapter, refitted arm by arm. Whatever they can and cannot do therefore
propagates directly into the effect estimate. Specifically, if either risk model
fails to capture an important pattern in the data — a region of the covariate
space where prediction is weak, or a source of variation left unexplained —
then subtracting two imperfect surfaces yields an effect estimate whose
heterogeneity cannot be trusted in that same region. Conversely, if both models
are strong and calibrated, their difference isolates the treatment effect
cleanly. This is why the quality of the effect estimates depends directly on
the accuracy of the outcome models. A T-learner cannot be sharper than the
two surfaces it differences. The second use is as a benchmark. The
causal forest, the Bayesian causal forest, the in-context model and the
interaction forest all fit outcome surfaces of their own, and the numbers in
this chapter are the reference against which their nuisance components are
judged.

\paragraph{One detail about what these models predict}
In the full-cohort models the randomized assignment is among the covariates,
so what is being predicted is the risk of an event \emph{under the treatment
the patient actually received}. That is the quantity Chapter 2 warns against
using as a decision rule, and it is the right quantity here: this chapter
measures how well outcomes can be anticipated, not how treatment changes them.
The separation between the two is the subject of Chapter 6, and the per-arm
models of Section~\ref{sec:per-arm} are the bridge to it.

\section{Protocol}

\paragraph{Models}
Four model families are fitted, each wrapped so that the five endpoints are
handled as five independent binary problems sharing one preprocessing chain: a
logistic regression, a random forest, a gradient boosting ensemble, and a soft
voting ensemble that averages the predicted probabilities of the other three.
The families are deliberately ordinary. The question this chapter answers is
how much signal the covariates carry, not which learner extracts it best, and
a spread of four families makes it possible to separate a property of the data
from a property of one algorithm.

\paragraph{Pipeline Preprocessing}
All models are evaluated using a threshold-free 5-fold
stratified cross-validation with respect to the target data over the full cohort of 4,579 patients. Within each fold, missing
values are imputed using a k-nearest-neighbour imputer and features are robustly scaled,
with both preprocessing steps fitted exclusively on the training data to avoid data leakage.
Since event rates range from 0.8\% to 11.1\%, the default 0.5 classification threshold would almost always
predict the majority class, making confusion-matrix-based metrics uninformative. We there-fore assess
discrimination without fixing an operating threshold, using AUROC and and the area under the precision--recall curve (AUPR).
Each patient is scored exactly once by a model that did not see them during training; 
the out-of-fold probabilities are then pooled to obtain a single ROC and PR curve, while mean and
standard deviation across folds are also reported to quantify variability, which is particularly
substantial for the rarest endpoints.

\paragraph{Hyperparameters}
Each family is additionally tuned with a Tree-structured Parzen Estimator
(Optuna), maximising mean average precision across the five endpoints under
3-fold cross-validation on the training split. Compared with each family's
default hyperparameters, tuning yields a small mean AUPR gain: $+0.014$ for
logistic regression, $+0.026$ for the random forest and $+0.006$ for gradient
boosting, none of which changes any of the conclusions of this chapter. The
gain from tuning is smaller than the difference in AUPR between endpoints
reported below: which endpoint is being predicted affects discrimination more
than which hyperparameters are used.

\paragraph{Reading AUPR}
Unlike AUROC, whose chance level is 0.5 for every endpoint, the chance level
of AUPR is the event rate itself. Every AUPR
below is accompanied by its \emph{lift}, the ratio between the achieved value
and the event rate, which is what makes the five endpoints comparable on the
rare positive class.

\paragraph{Calibrating the per-arm models}
The per-arm outcome models built for the T-learner are additionally wrapped,
after tuning, in a Platt (sigmoid) calibrator \citep{platt1999}, fitted by internal
cross-validation on the training fold, rather than an isotonic one, because
isotonic fitting is unstable at these event rates.

\section{Discrimination}

Table~\ref{tab:risk-auroc} reports the out-of-fold AUROC of the four families
on the five endpoints, with the endpoints ordered by decreasing event count.
Figure~\ref{fig:roc-pr} shows the pooled curves for the soft voting ensemble,
which is the strongest of the four on average and is used as the reference
model in the rest of the chapter.

\begin{table}[h]
\centering
\caption{Out-of-fold AUROC (mean $\pm$ std across 5 folds), full cohort.
Endpoints ordered by decreasing event count.}
\label{tab:risk-auroc}
\begin{tabular}{lrcccc}
\toprule
Endpoint & Events & LR & RF & GB & Soft vote \\
\midrule
Bleeding             & 506 & $0.62 \pm 0.03$ & $0.58 \pm 0.02$ & $0.61 \pm 0.04$ & $0.62 \pm 0.03$ \\
Bleeding BARC 2/3/5  & 359 & $0.62 \pm 0.03$ & $0.57 \pm 0.04$ & $0.59 \pm 0.02$ & $0.63 \pm 0.02$ \\
MI                   & 109 & $0.72 \pm 0.05$ & $0.71 \pm 0.04$ & $0.70 \pm 0.04$ & $0.73 \pm 0.04$ \\
Cardiovascular death &  81 & $0.76 \pm 0.05$ & $0.69 \pm 0.05$ & $0.73 \pm 0.05$ & $0.76 \pm 0.04$ \\
Stroke               &  35 & $0.60 \pm 0.08$ & $0.62 \pm 0.11$ & $0.67 \pm 0.13$ & $0.60 \pm 0.03$ \\
\bottomrule
\end{tabular}
\end{table}

\begin{figure}[h]
\centering
\includegraphics[width=1\textwidth]{ROC_PR_out_of_fold.png}
\caption{Pooled out-of-fold receiver operating characteristic (left panel)
and precision-recall (right panel) curves for the soft voting ensemble,
all five endpoints. Each endpoint is identified by its name from
Table~\ref{tab:risk-auroc}: bleeding (light blue), bleeding BARC 2/3/5 (dark blue),
myocardial infarction (dark red), cardiovascular death (green), and stroke
(grey). Every patient is scored once by a model fitted without them. In
the right-hand panel the dotted horizontal lines mark the chance level of each
curve, which is the endpoint's event rate.}
\label{fig:roc-pr}
\end{figure}

The models carry a genuine but limited signal for every endpoint: all five curves  in the ROC
panel of Figure~\ref{fig:roc-pr} sit above the diagonal, and all five AUPR
values sit above their event rate (shown by the dotted lines in the precision-recall
panel). However, the discrimination achieved is weak to moderate across all
endpoints. The covariates are informative, but their predictive capacity is limited.


\begin{table}[h]
\centering
\caption{Out-of-fold AUPR for the soft voting ensemble, with its event rate and the resulting lift.}
\label{tab:risk-auprc}
\begin{tabular}{lccc}
\toprule
Endpoint & AUPR & event rate & Lift \\
\midrule
Bleeding             & 0.175 & 0.111 & 1.58 \\
Bleeding BARC 2/3/5  & 0.127 & 0.078 & 1.62 \\
MI                   & 0.073 & 0.024 & 3.07 \\
Cardiovascular death & 0.059 & 0.018 & 3.33 \\
Stroke               & 0.028 & 0.008 & 3.70 \\
\bottomrule
\end{tabular}
\end{table}

\section{An Anomaly: More Events, Worse Prediction}

Table~\ref{tab:risk-auroc} shows an unexpected ordering: endpoints with fewer events are predicted better than bleeding, 
despite having substantially less positive samples. Cardiovascular death (81 events) reaches an AUROC of 0.76 
and myocardial infarction (109 events) 0.73, compared with 0.62 for bleeding (506) and 0.63 for BARC 2/3/5 (359). 
The same pattern appears in the AUPR lift, where the rarer endpoints show
stronger enrichment relative to their chance levels.

This contradicts the simple expectation that more events should provide more 
information and therefore better prediction. It is also clinically unexpected, 
since bleeding is the main endpoint of the trial and the covariates include several 
variables directly related to bleeding risk.

\paragraph{It is not an artefact of the model family.}
The same ordering is reproduced by all four model families in Table~\ref{tab:risk-auroc}. 
The AUROC gap between cardiovascular death and bleeding ranges from 0.11 to 0.14 across models. 
This is consistent with an independent earlier analysis \citep{maroni2026} using gradient-boosted trees and a different validation scheme
found the same pattern, with myocardial infarction predicted better than both bleeding definitions.
The absolute values differ, but the ordering remains.

\paragraph{Stroke is too unstable to interpret.}
Stroke, with only 35 events, provides no reliable evidence. With roughly seven 
events per fold 35 events in total, its AUROC varies substantially across folds, with standard 
deviations between 0.03 and 0.13 and means between 0.60 and 0.67. 
The estimate is therefore too unstable to support a meaningful conclusion.
If we observe the Figure \ref{fig:roc-pr} curves, the stroke ROC curve is very
unstable and almost overlaps with the diagonal. With this many
positive events is really hard to get a stable estimate of the model's performance.

\paragraph{What the models look at.}
The average absolute SHAP feature attributions \citep{lundberg2017shap} computed on logistic regression, random forest and gradient boosting
 are consistent with the anomaly presented 
at the beginning of this section rather than
explaining it away. For cardiovascular death (Figure~\ref{fig:shap-cvdeath})
the leading contributors are continuous physiological measurements —
haemoglobin, ejection fraction, creatinine, heart rate, contrast volume — plus
recent myocardial infarction, all of which vary widely across the cohort. For
bleeding (Figure~\ref{fig:shap-bleeding}) the leading contributors are age,
oral anticoagulation, renal function and the
acute-coronary-syndrome presentation; several of these variables overlap with
the eligibility criteria through which the cohort was assembled \citep{valgimigli2021appendix}.

\begin{figure}[h]
\centering
\includegraphics[width=1\textwidth]{SHAP_cvdeath.png}
\caption{Mean absolute SHAP attribution, top ten covariates, cardiovascular
death, one panel each for logistic regression, random forest and gradient
boosting. The ranking is dominated by continuous physiological markers.}
\label{fig:shap-cvdeath}
\end{figure}

\begin{figure}[h]
\centering
\includegraphics[width=1\textwidth]{SHAP_bleeding.png}
\caption{Mean absolute SHAP attribution, top ten covariates, bleeding, one
panel each for logistic regression, random forest and gradient boosting. The
ranking includes variables that also appear among the trial's
high-bleeding-risk eligibility criteria \citep{valgimigli2021appendix}.}
\label{fig:shap-bleeding}
\end{figure}

\section{Per-arm Models for the T-learner}
\label{sec:per-arm}

The models above are fitted on the whole cohort. The two-model meta-learner of
Chapter 6 needs something slightly different: one outcome model per treatment
arm, so that their difference estimates the effect of the treatment. Chapter 2
sets out what this requires of the two arm-specific models: the same model
family on both arms, and calibrated probabilities. Both are met here, by
construction, since they are properties of how the risk models of this
chapter are fitted rather than of the effect estimator applied to them
afterwards.

\paragraph{The same family on both arms.}
Both arms are fitted with a single model family, tuned separately on each
arm's training data.

\paragraph{Calibrated probabilities.}
Each arm's tuned pipeline is calibrated as described in the protocol above.
The out-of-fold Brier score improves substantially on every endpoint and in
both arms, for instance from 0.189 to 0.118 for bleeding under prolonged
therapy and from 0.167 to 0.081 under abbreviated therapy.

\paragraph{Discrimination arm by arm.}
Table~\ref{tab:per-arm} reports the out-of-fold AUROC within each arm,
together with the arm's raw event rate. Splitting the cohort halves the data
available per model, so the estimates are noisier than those of
Table~\ref{tab:risk-auroc}; what survives the split is the ordering. Within
both arms, cardiovascular death and MI are predicted better than either
bleeding endpoint, so the anomaly is a property of the cohort and not of
pooling the two arms.

\begin{table}[h]
\centering
\caption{Per-arm event rate and out-of-fold AUROC (mean $\pm$ std across
5 folds) for the T-learner's component models, 2,284 patients under prolonged
therapy and 2,295 under abbreviated therapy.}
\label{tab:per-arm}
\begin{tabular}{lcccc}
\toprule
 & \multicolumn{2}{c}{Prolonged ($T=1$)} & \multicolumn{2}{c}{Abbreviated ($T=0$)} \\
\cmidrule(lr){2-3}\cmidrule(lr){4-5}
Endpoint & Event rate & AUROC & Event rate & AUROC \\
\midrule
Bleeding             & 13.3\% & $0.59 \pm 0.02$ &  8.8\% & $0.61 \pm 0.05$ \\
Bleeding BARC 2/3/5  &  9.2\% & $0.59 \pm 0.04$ &  6.4\% & $0.62 \pm 0.06$ \\
MI                   &  2.1\% & $0.73 \pm 0.04$ &  2.6\% & $0.69 \pm 0.04$ \\
Cardiovascular death &  1.9\% & $0.76 \pm 0.08$ &  1.6\% & $0.73 \pm 0.07$ \\
Stroke               &  1.0\% & $0.73 \pm 0.03$ &  0.5\% & $0.56 \pm 0.18$ \\
\bottomrule
\end{tabular}
\end{table}

The event-rate columns of Table~\ref{tab:per-arm} are worth a moment on their
own, because they are the trial's own result stated without any model. Subtracting the abbreviated-arm 
event rate from the prolonged-arm event rate gives an
unadjusted between-arm difference of $+4.5$ percentage points on overall bleeding and $+2.8$ on BARC 2/3/5 indicating a higher bleeding incidence
under prolonged therapy. For ischemic targets the the corresponding differences are $-0.5$, $+0.3$ and $+0.5$
percentage points on MI, cardiovascular death and stroke, respectively. These differences
are small relative to the observed event counts and do not provide evidence
of a detectable treatment effect, consistent with the conclusions of the
original trial report \citep{valgimigli2021masterdapt}.

\section{The Anomaly, Left Open}

This chapter has produced the components the next one needs, and one finding
it cannot explain on its own. Bleeding, the endpoint with the most events,
the one the trial was designed around, and the one whose predictors are best
represented in the covariate set, is the endpoint predicted worst, in every
model family, in both treatment arms, and in an independent earlier modelling
pass.
Cardiovascular death in is predicted best even if the provided positive
samples are much
fewer than the one present for the bleeding endpoints.

Three readings are available at this point. The covariate set might simply
lack the variables that determine bleeding, which the feature attributions
make implausible: the models are using precisely the variables a clinician
would name. The bleeding endpoint might be intrinsically less predictable than
death, which is not what the PRECISE-DAPT derivation and validation cohorts
report: its C-index for out-of-hospital bleeding ranged from 0.66 to 0.73
\citep{costa2017precisedapt}. Or the cohort itself might have been assembled in
a way that removes the variation a bleeding model would need.

\section{Chapter Summary}
This chapter explores the results of the risk models fitted on the full cohort.
Are described: the preprocessing steps of the pipeline, how the models are tuned,
calibrated and evaluated.
The results show that the models carry a limited but present signal for every endpoint
with AUROC values between 0.60 and 0.76 and AUPR lifts between 1.58 and 3.70.
In general the ischemic endpoints are predicter with better precision than 
the bleeding endpoints, despite having fewer events.
This is an unexpected result, since the bleeding endpoints should be the primary focus of the trial.
Is included a discussion on the feature importance of the models used to predict the endpoints using SHAP values.
Finally, the per-arm models used in the T-learner are described, 
showing that the same anomaly is present when the cohort is split by treatment arm.